\FloatBarrier

\section{Introduction}

Carpal tunnel syndrome affects between 1 \% and 5 \% of the general adult population \CITE{cts} and a substantially larger fraction of intensive computer users, and upper-limb motor impairment arising from spinal cord injury, amyotrophic lateral sclerosis, cerebral palsy or stroke removes conventional pointing devices from reach altogether. For both populations the operative requirement is the same: the full command surface of a personal computer --- cursor positioning, clicking, application and window switching, text entry, scrolling --- must be reachable without the hands. Systems that meet this requirement by substituting \emph{gaze} for the cursor, \emph{blink} for the click, \emph{head pose} for discrete gestures and \emph{voice} for command entry have been demonstrated repeatedly in the literature, yet deployment outside the laboratory remains rare. The obstacles are not principally in the individual perception components, each of which has strong published solutions; they lie in the joins between them, and in how the evidence for them is produced.

Three obstacles recur across the literature and are the organising concerns of this work.

\begin{itemize}
\item \textbf{Evaluation that does not survive an unseen user.} Appearance-based perception corpora are group-structured: the MRL Eye Dataset \CITE{mrl} holds roughly 2,300 frames of the same eye per subject, MPIIGaze \CITE{mpiigaze,mpiigaze_pami} several thousand normalised patches per participant, and 300W-LP \CITE{w300lp} synthesises a median of 34 viewpoint renders of each of its 3,837 base identities. Partitioning such corpora by random row sampling places near-duplicate observations of one individual on both sides of the split, reducing the test task from generalisation to nearest-neighbour retrieval against an almost identical stored exemplar. The remedy --- group-disjoint partitioning keyed on subject, participant or face identity --- is standard and well documented \CITE{leakage,kapoor}. What is not standard, and what this paper argues is the more consequential omission, is that applying the remedy correctly to \emph{every model in a comparison} does not make the comparison correct. When a proposed model and its pretrained baseline were developed on overlapping corpora, the evaluation sample can lie inside one model's training partition and outside the other's, and the resulting contrast measures exposure asymmetry rather than capability. Nothing in the per-pipeline audits prescribed by the standard taxonomies \CITE{kapoor,leaktypes} detects this, because each such audit sees only one pipeline and neither pipeline carries any representation of the other's partition boundary.

\item \textbf{A budget that must accommodate concurrency.} A hands-free interface is not one perception problem but four, running simultaneously, on the machine the user already owns, without an accelerator. The camera loop at 30 frames per second affords a 33 ms budget within which face landmarking, three vision networks, two finite-state machines and a cursor stabilisation chain must all complete; a fourth modality runs on a separate audio thread. Accuracy obtained at a cost that forecloses concurrency is not usable accuracy, so the operative question is not which model is most accurate but which accuracy--footprint operating points admit four simultaneous modalities.

\item \textbf{Component metrics that do not predict composite behaviour.} No perception network in such a system is read directly. Each feeds a decision rule with its own thresholds, dwell times, hysteresis bands and filters, and it is that composite the user experiences. A per-frame blink classification is not a click; a click is a closure sustained past 420 ms, roughly thirteen consecutive frames, with natural blinks of 100--300 ms rejected by construction. A head-pose estimate is not a gesture; a gesture is a threshold crossing held for 650 ms inside an 18$^{\circ}$ hysteresis band. A raw gaze angle is not a cursor position; it is the input to a per-user polynomial calibration corrected online by recursive least squares \CITE{rls} and smoothed by a one-euro filter \CITE{onceuro}. Reporting the component metric and reasoning about the composite as though the two were interchangeable inverts several conclusions, as Section 5 demonstrates in three separate places.
\end{itemize}

These three obstacles are treated in isolation in the existing literature. Work that strengthens evaluation protocol rarely reports deployment cost; work that optimises for edge inference rarely audits its splits; work that adds interpretability rarely connects the explanation to the decision rule that consumes the model's output. The fragmentation, rather than any single gap, is what keeps demonstrated hands-free systems in the laboratory.

\FloatBarrier

\subsection{Contributions}

This paper presents MAPLE (Multimodal Audited Perception with Leakage-aware Evaluation), a framework that treats the three obstacles as one problem, and reports the evaluation of a four-modality hands-free interaction system conducted under it. The contributions are as follows.

\begin{enumerate}
\item \textbf{Formalisation of asymmetric exposure as a distinct leakage class.} We define a cross-model comparison as \emph{asymmetrically exposed} when three conditions hold simultaneously: the compared models were developed on overlapping corpora; the corpus carries group structure with strong within-group dependence; and the evaluation sample is drawn from groups lying on both sides of at least one compared model's training-partition boundary. We prove (Theorem 1) that under these conditions the affected model's aggregate score is a convex combination of a memorisation regime and a generalisation regime whose mixing weight is fixed by a sampling quota, so that the reported accuracy can be changed without changing any model, any datum or any metric definition. We show that neither of the two standard leakage taxonomies \CITE{kapoor,leaktypes} contains a category for this fault, and that the existing remedies --- lineage tracing \CITE{tsfm} and contamination search \CITE{oren} --- presume documentation or corpus access an external auditor rarely has.

\item \textbf{GERA: a prediction-only detector with a negative control.} We introduce the Group Error-Ratio Audit, a statistic \emph{r}(\emph{g}) formed as the ratio of an audited model's per-group mean error to a reference model's mean error on the same group. We prove (Proposition 2) that under a multiplicative difficulty model the intrinsic difficulty term cancels exactly, leaving a quantity that varies across groups only through differential familiarity. GERA adapts difficulty calibration --- developed for strengthening membership inference attacks \CITE{shokri,watson} --- from example level to group level, which is the shift that makes it stable enough to read off published predictions. It requires only the two prediction streams over a shared evaluation set and a group label per sample: no training data, no lineage documentation, no corpus access. On the system studied here it separates training-partition membership at AUC 0.889 with Mann--Whitney \emph{p} = 0.006, and a negative control on a modality where neither model saw the evaluation groups returns \emph{I}$^{2}$ = 0.0 \% against 99.8 \% for the contaminated comparison, establishing that it stays silent when it should even under a fourfold swing in intrinsic difficulty.

\item \textbf{Decomposition of the exposure bias into difficulty and familiarity.} We show (Theorem 3) that the same reference model that calibrates the detector also decomposes the observed inflation into an intrinsic-difficulty and a familiarity component. The split depends on a transfer coefficient the data do not identify, so we report it as an interval rather than a point estimate: of a 1.905$^{\circ}$ aggregate inflation on the affected modality, the familiarity component is 0.425$^{\circ}$--1.068$^{\circ}$ (22 \%--56 \%), the ends corresponding to a multiplicative and an additive difficulty model respectively. Under both, intrinsic difficulty fails to account for the whole gap. This is a quantity lineage analysis cannot supply at all, since knowing \emph{that} a model saw a corpus does not say how much of its advantage came from having seen it.

\item \textbf{A six-step corrected comparison protocol (EDCP), and its exercise on four modalities.} We state an Exposure-Disjoint Comparison Protocol whose first two steps are established practice, whose steps three through five formalise for group-structured corpora the discipline that lineage analysis prescribes informally, and whose sixth step adds the prediction-only audit. We then exercise the protocol on four perception modalities spanning two regression and two classification tasks, reporting for each the published score under the faulty protocol, a fair baseline retraining the original hyperparameters under the corrected protocol, and a tuned score under that same protocol, so that the protocol effect and the optimisation effect can never be conflated.

\item \textbf{Four compact task-specific networks characterised at their deployed operating point.} We design, tune and evaluate four convolutional networks trained from scratch --- not fine-tuned backbones --- spanning 101,766 to 1,142,339 parameters and 7.84 MB in aggregate, each deployed as a self-contained TorchScript module \CITE{torchscript,pytorch} running single-threaded on CPU, so that every latency and accuracy figure is measured on the artefact that ships rather than on a research-time proxy. Aggregate vision-pipeline cost is 6.68 ms per frame inside a 33 ms budget, which is the condition that permits four modalities to run concurrently without an accelerator.

\item \textbf{A unified mechanistic account of the residual error.} We show that a single parameter --- the fitted regression gain of predicted on true target --- governs the error structure of both regression modalities, accounting for the isotropic radial growth of gaze error, the constant-ratio degradation of head-pose error across the pose manifold, the axis-ranking inversions between absolute and variance-normalised metrics, and the absence of any localised failure mode. We further show that inverting the fitted gain \emph{increases} angular error by 28.5 \%, establishing that the contraction is the estimator's variance-optimal response to representational uncertainty rather than a defect, and thereby redirecting engineering effort from calibration to representation. For the two classification modalities we show the opposite diagnosis: ranking quality far exceeds operating-point accuracy, placing their residual error in decision logic, which is why a confidence gate and a threshold move improve them where additional capacity would not.

\item \textbf{Decision-rule-aware reporting, which inverts three conclusions.} By stating for each modality the deployed decision rule and the error statistic that rule makes diagnostic, we obtain three conclusions the component metrics alone would contradict: head pose survives an 8.74$^{\circ}$ yaw error because its error is unbiased and its gesture rule averages roughly twenty frames inside an 18$^{\circ}$ hysteresis band; blink inverts the conventional cost asymmetry, so thresholds \emph{below} the default are defensible because an isolated false positive cannot fire a click whereas a false negative breaks a deliberate closure run; and the weakest keyword in the voice vocabulary is bound to the only command that reactivates a paused interface, making the largest error mode a vocabulary-design fault removable at zero computational cost.

\item \textbf{A demonstration that the protocol fault is consequential rather than pedantic.} Under the uncorrected protocol the proposed gaze network scored 2.84$^{\circ}$, level with its pretrained baseline's all-participant figure of 2.835$^{\circ}$, and the decision was to retain it. Under the corrected protocol the true held-out contrast is 6.439$^{\circ}$ against 3.621$^{\circ}$, the baseline is simultaneously 7.1$\times$ smaller, and the recommendation reverses to adopting the pretrained model. We know of no sharper available demonstration that the distinction between a per-pipeline audit and a per-comparison audit changes engineering outcomes.
\end{enumerate}

\FloatBarrier

\subsection{Organisation}

Section 2 reviews related work along the three dimensions corresponding to the obstacles above --- hands-free interaction modalities and their cross-subject behaviour, evaluation integrity and the leakage literature, and efficient and interpretable on-device inference --- closing with a comparative synthesis that identifies the unresolved gaps. Section 3 develops the MAPLE methodology: the system formalisation, the datasets and their grouping keys, the preprocessing and input representations, group-disjoint partitioning, the formal treatment of asymmetric exposure including the GERA statistic and the decomposition theorem, the four network architectures, the training objectives and search procedure, the deployed decision rules, the attribution method, the evaluation protocols and metric definitions, and a complexity analysis. Section 4 states the experimental settings, implementation environment, baselines and statistical procedures. Section 5 reports and discusses the results across all four modalities, the joint findings, and the accuracy--efficiency operating point. Section 6 states threats to validity, and Section 7 concludes with limitations and directions for future work.
